% !PRINTED 474
\weberpage{474}
%<<CONT
\gesperrt{Kronecker} erkannt\footnote{\gesperrt{Kronecker}, \glqq Ueber Gleichungen 7\textsuperscript{ten} Grades\grqq{}. Monatsberichte der Berliner Akademie 1858.}, und seine nähere Untersuchung ist für die allgemeine Theorie der Gleichung 7\textsuperscript{ten} Grades von grosser Wichtigkeit.

Wir können diese Permutationsgruppe dadurch erhalten, dass wir die Permutationen aufsuchen, die durch die Substitutionen der Gruppe $G_{168}$ unter den Grössen $z_0, z_1, z_2, z_3, z_4, z_5, z_6$ (\S.\,125) hervorgerufen werden. Hierbei ist dann in Bezug auf die Zusammensetzung zu beachten, dass, wenn die beiden linearen Substitutionen $\xi_1, \xi_2$ die Permutationen $\pi_1, \pi_2$ bewirken, die zusammengesetzte Permutation $\pi_1 \pi_2$ durch $\xi_2 \xi_1$ hervorgerufen wird. Denn nach der Definition erhält man $\xi_2 \xi_1 (x)$ dadurch, dass man auf die Variablen $(x)$ zunächst die Substitution $\xi_1$ und auf das Ergebniss die Substitution $\xi_2$ anwendet. Ebenso bedeutet $\pi_1 \pi_2$ die Permutation, die sich ergiebt, wenn man auf die sieben Ziffern zuerst $\pi_1$ und darauf $\pi_2$ anwendet. Wir erhalten so, der Gruppe $G_{168}$ entsprechend, eine Permutationsgruppe 168\textsuperscript{sten} Grades, die wir mit $P_{168}$ bezeichnen, und diese beiden Gruppen sind isomorph.

Da wir $\tau$ und $\omega$ als erzeugende Elemente der Gruppe $G_{168}$ erkannt haben (\S.\,72, I.), so genügt es, wenn wir die diesen beiden Substitutionen entsprechenden Permutationen bestimmen, um daraus die ganze Gruppe $P_{168}$ abzuleiten.

Nun geht aber aus der Substitution $\tau$ die cyklische Permutation der sieben Ziffern

\[
(\tau) \qquad (0, 1, 2, 3, 4, 5, 6)
\]

hervor, und der Einfluss von $\omega$ ergiebt sich aus der Bemerkung, dass $z_0$ durch die Substitutionen der Octa\"edergruppe $\chi^\lambda \omega^u \Theta^v$ ungeändert bleibt. Um also die Aenderung von $z_r$ durch $\omega$ zu erhalten, haben wir nur den Einfluss von $\omega \tau^r$ auf $z_0$ zu ermitteln. Dieser Einfluss aber ergiebt sich unmittelbar aus den Formeln \S.\,72, (15), wonach z. B. $\omega \tau = \tau^2 \chi^2 \omega \Theta^2$ ist, so dass also $z_1$ in $z_2$ übergeht u. s. f. Demnach entspricht der Substitution $\omega$ die Permutation

\[
(\omega) \qquad \begin{pmatrix} 0, & 1, & 2, & 3, & 4, & 5, & 6 \\ 0, & 2, & 1, & 5, & 4, & 3, & 6 \end{pmatrix} = (1, 2)\,(3, 5).
\]

Man kann die Gruppe $P_{168}$ durch die Congruenzgruppe ter-
%CONT>>
